\documentclass[11pt,a4paper]{article}
\usepackage[utf8]{inputenc}
\usepackage[T1]{fontenc}
\usepackage[left=1.25cm, top=1.2cm, right=1.25cm, bottom=1.2cm]{geometry}
\usepackage{enumitem}
\usepackage{hyperref}
\usepackage{titlesec}
\usepackage{xcolor}

% --- Tipografía Corporativa y Moderna ---
\usepackage{mathpazo} 
\usepackage[scaled=0.95]{helvet}

% --- Style Configuration ---
\titleformat{\section}{\large\bfseries\uppercase}{}{0pt}{}[\titlerule]
\titlespacing{\section}{0pt}{16pt}{8pt}

% --- Control de Espaciados Automatizado ---
\setlength{\parskip}{4pt}
\setlist[itemize]{leftmargin=1.3em, itemsep=2pt, parsep=0pt, topsep=2pt, partopsep=0pt}

% Comando personalizado para separar las experiencias laborales de forma limpia
\newcommand{\companygap}{\vspace{12pt}}

\begin{document}

% --- HEADER ---
\begin{center}
    {\Huge \textbf{Rigoberto Barra}} \\
    \vspace{3pt}
    {\large \textbf{Analytics Engineer \textbar{} Data Engineer}} \\
    \vspace{3pt}
    Viña del Mar, Chile \textbar{} +56-996974170 \textbar{} \href{mailto:rigbarra@outlook.com}{rigbarra@outlook.com} \textbar{} \href{https://www.linkedin.com/in/rigbarra}{linkedin.com/in/rigbarra}
\end{center}

% --- PROFESSIONAL SUMMARY ---
\section*{Resumen Profesional}
Ingeniero Industrial con más de 10 años de experiencia actuando como puente estratégico entre unidades técnicas y de negocio. Especializado en diseñar arquitecturas de datos escalables y productos analíticos de alto impacto utilizando el Modern Data Stack (SQL, Python, AWS Serverless, dbt, Power BI). Experiencia comprobada alineando la ejecución técnica con la estrategia del negocio mediante el liderazgo de migraciones a la nube y proyectos de automatización a gran escala. Poseo Diplomados en Data Science (PUC) y Business Intelligence (U. de Chile). Inglés Profesional (B2+) consolidado tras 2 años de experiencia viviendo y trabajando en Dublín, Irlanda.

% --- TECHNICAL SKILLS ---
\section*{Habilidades Técnicas}
\begin{itemize}
    \item \textbf{Data Engineering \& Modeling:} SQL (Joins complejos, CTEs, Window Functions), Python (Pandas, Polars), DuckDB, dbt, Dagster, Docker.
    \item \textbf{Cloud \& Big Data Stack:} AWS (Lambda, Step Functions, Glue, Athena, S3, Redshift, IAM), GCP (BigQuery).
    \item \textbf{BI, Analytics \& Visualization:} Power BI (DAX), Apache Superset, Quarto, Figma (UX Layouts).
    \item \textbf{Data Quality \& Governance:} Git (GitHub), Testing (Pytest), CI/CD Pipelines (GitHub Actions).
\end{itemize}

% --- PROFESSIONAL EXPERIENCE ---
\section*{Experiencia Profesional}

\noindent
\textbf{Data Engineer Specialist} \hfill Mar 2026 – Jun 2026 \\
\small \textit{Banco Estado \textbar{} Contrato por Proyecto / Plazo Fijo} \hfill \small Santiago, Chile
\begin{itemize}
    \item \textbf{Diseñé e implementé} pipelines ETL serverless utilizando AWS (Lambda, Glue, Step Functions, Athena) bajo un esquema Medallion (capas Bronze, Silver, Gold) para automatizar la información financiera crítica del departamento de Contabilidad, integrando un Data Lake interno y APIs externas.
    \item \textbf{Lideré el diseño y modelamiento del ecosistema analítico}, abarcando la maquetación de la experiencia de usuario (UX) en Figma y el desarrollo de tableros avanzados y modelos semánticos en Power BI para los gestores contables.
    \item \textbf{Garanticé la robustez y calidad de los datos} mediante la implementación de suites de pruebas automatizadas con PyTest para verificar y certificar la lógica de las transformaciones antes del despliegue del código.
    \item \textbf{Aceleré la velocidad de desarrollo} y la refactorización de código mediante flujos de trabajo asistidos por IA (LLMs), estableciendo una catalogación de datos exhaustiva y documentación técnica generada a través de Quarto (HTML).
\end{itemize}

\companygap
\noindent
\textbf{Data Engineer} \hfill Ene 2025 – Sep 2025 \\
\small \textit{Fortis Abroad (Consultoría Intercambio Estudiantil)} \hfill \small Dublín, Irlanda
\begin{itemize}
    \item \textbf{Diseñé e implementé} la arquitectura de datos relacional local utilizando DuckDB y dbt para centralizar, transformar y modelar flujos globales fragmentados de información académica multi-región, estructuras de precios y conversión de divisas en una Única Fuente de Verdad (SSOT).
    \item \textbf{Orquesté} pipelines end-to-end mediante Dagster a nivel local, automatizando la ingesta diaria de datos operacionales desde APIs externas, el manejo de dependencias de archivos y la ejecución de scripts de transformación en Python.
    \item \textbf{Diseñé y automatizé} un cotizador interactivo y reportes operacionales mediante Quarto, generando aplicaciones HTML dinámicas embebidas que redujeron drásticamente los tiempos de respuesta del equipo de ventas, operando al 100\% en un entorno en inglés.
\end{itemize}

\companygap
\noindent
\textbf{Analytics Engineer} \hfill Nov 2022 – Jul 2023 \\
\small \textit{CHEK (Banco Ripley)} \hfill \small Santiago, Chile
\begin{itemize}
    \item \textbf{Lideré la migración} y remodelación de flujos de datos legacy desde SAS hacia AWS Redshift, habilitando capacidades analíticas en tiempo real a través de la empresa.
    \item \textbf{Diseñé y desplegué} una arquitectura integral de Customer 360 (+100 atributos), acelerando la entrega de KPIs ejecutivos para la toma de decisiones estratégicas en un 50\%.
    \item \textbf{Impulsé la adopción de la billetera digital en un 15\%} e incrementé el ROI en un 10\% mediante modelos de segmentación de clientes basados en Python y rigurosos tests A/B.
    \item \textbf{Fomenté la colaboración cross-functional} con equipos comerciales para co-diseñar estrategias de activación de usuarios y campañas de marketing personalizadas basadas en analítica.
\end{itemize}

\companygap
\noindent
\textbf{Customer Intelligence Manager} \hfill Jun 2021 – Sep 2022 \\
\small \textit{SMU (Unimarc)} \hfill \small Santiago, Chile
\begin{itemize}
    \item \textbf{Orquesté la unificación de datos transaccionales} mediante pipelines en Python y R, estableciendo con éxito una Única Fuente de Verdad (SSOT) para toda la organización. 
    \item \textbf{Desarrollé e implementé} un framework de valor incremental (Lift) para clientes de la tarjeta Unipay, impulsando un aumento del 40\% en el gasto promedio de los grupos segmentados vs. control.
    \item \textbf{Lideré el diseño técnico} de metodologías de A/B Testing, aumentando el engagement de usuarios en un 12\% al alinear la arquitectura de datos con los objetivos comerciales de conversión.
    \item \textbf{Garanticé la calidad de la entrega técnica} y desarrollé indicadores clave de comportamiento para stakeholders ejecutivos de la división de crédito del grupo.
\end{itemize}

\companygap
\noindent
\textbf{Senior Churn Analyst} \hfill Dic 2017 – May 2021 \\
\small \textit{DIRECTV (AT\&T)} \hfill \small Santiago, Chile
\begin{itemize}
    \item \textbf{Construí un modelo de Payment Behavior Scoring} que automatizó estrategias de "no intervención", logrando una reducción del 40\% en la carga de trabajo del call center.
    \item \textbf{Definí estándares analíticos} para la región del Cono Sur (CL, AR, PE, UY) bajo marcos de trabajo Agile, escalando modelos de propensión y reportería iterativa a través de los mercados.
\end{itemize}

\companygap
\noindent
\textbf{Business Intelligence Analyst} \hfill Mar 2015 – Dic 2017 \\
\small \textit{VTR (Liberty Global)} \hfill \small Santiago, Chile
\begin{itemize}
    \item \textbf{Desarrollé un ecosistema de datos GIS} en Power BI para impulsar estrategias comerciales y operacionales basadas en KPIs geoespaciales.
    \item \textbf{Audité procesos técnicos} mediante SQL y SPSS Modeler, eliminando con éxito 12.000 llamadas mensuales redundantes de clientes.
\end{itemize}

% --- EDUCATION ---
\section*{Educación}

\noindent
\textbf{Curso de Inglés General} \hfill 2023 – 2025 \\
\small ISI Dublin English School \hfill \small Dublín, Irlanda

\vspace{6pt}

\noindent
\textbf{Diplomado en Data Science} \hfill 2021 \\
\small Pontificia Universidad Católica de Chile \hfill \small Santiago, Chile

\vspace{6pt}

\noindent
\textbf{Diplomado en Business Intelligence} \hfill 2018 \\
\small Universidad de Chile \hfill \small Santiago, Chile

\vspace{6pt}

\noindent
\textbf{Ingeniería Civil Industrial} \hfill 2008 – 2014 \\
\small Universidad del Bío-Bío \hfill \small Concepción, Chile

\end{document}